\subsection{Whole original analytic points on arbitrary homogeneous charts}

\begin{definition}[The original ambient fraction homomorphism]
\label{projective:actual-admissible-ambient-fraction-hom}
\lean{EnvelopeProperty.ProjSubvariety.nativeAdmissibleAmbientChartHom}
\leanok
\uses{projective:actual-original-admissible-chart-scheme-iso}
Let $X$ be the original projective subvariety over a field, and let $F$ be
homogeneous of degree $d$, nonzero modulo its original homogeneous ideal.
There is a canonical ring homomorphism from the original ambient homogeneous
localization $P_{(F)}$ to the original quotient-chart algebra $A_{X,F}$.
It is the actual graded quotient map followed by the inverse of the actual
quotient-chart/native-away algebra equivalence. No normality assumption is used.
\end{definition}
\begin{proof}
Apply homogeneous localization to the original graded quotient map, then
compose with the inverse ring homomorphism of the established whole algebra
equivalence. Both maps retain the original coefficient actions.
\end{proof}

\begin{lemma}[Every original numerator is retained]
\label{projective:actual-admissible-ambient-fraction-numerator}
\lean{EnvelopeProperty.ProjSubvariety.nativeAdmissibleAmbientChartHom_mk}
\leanok
\uses{projective:actual-admissible-ambient-fraction-hom,projective:actual-admissible-native-homogeneous-fraction}
For every $m\in\mathbb N$ and every original homogeneous $G$ of degree $md$,
the ambient fraction $G/F^m$ maps to precisely the original quotient-chart
class $\operatorname{quotChartMk}(m,G)$. This is equality of ring elements.
\end{lemma}
\begin{proof}
The actual homogeneous localization quotient map sends the fraction to the
fraction of the original quotient images. Apply injectivity of the established
whole quotient-chart equivalence and its exact numerator formula. Cancel the
inverse equivalence with the equivalence itself.
\end{proof}

\begin{lemma}[The whole original inclusion square]
\label{projective:actual-admissible-ambient-native-square}
\lean{EnvelopeProperty.ProjSubvariety.nativeAdmissibleChartMap_inclusion}
\leanok
\uses{projective:actual-admissible-ambient-fraction-hom,projective:actual-original-admissible-chart-scheme-iso}
For $d>0$, the actual map $\operatorname{Spec} A_{X,F}\to X$ followed by the
original closed projective inclusion equals the spectrum of the preceding
ambient fraction homomorphism followed by the actual open immersion
$\operatorname{Spec} P_{(F)}\to\operatorname{Proj} P$.
\end{lemma}
\begin{proof}
Unfold the actual chart map through the quotient Proj comparison. Its composite
with the original inclusion is the genuine graded quotient Proj map. Use the
native homogeneous-localization square for that map. Functoriality of spectrum
identifies the composite of the two affine maps with the spectrum of their
literal composed ring homomorphism, which is the preceding definition.
\end{proof}

\begin{theorem}[The original analytic chart inverse is the actual whole native point]
\label{projective:actual-admissible-analytic-native-chart}
\lean{EnvelopeProperty.ProjSubvariety.nativeAnalyticHomeomorph_chartInv}
\leanok
\uses{projective:actual-admissible-ambient-fraction-numerator,projective:actual-admissible-ambient-native-square,projective:actual-graded-native-canonical-away-seminorm}
Let $z$ be any original affine seminorm point of $A_{X,F}$, with $F$ as above
and $d>0$. Under the original projective analytic/native homeomorphism, the
original chart inverse of $z$ is exactly the whole native point obtained by
mapping the actual affine native point of $z$ through the original chart map.
In particular, its original support and residue-field valuation are retained;
no point-preservation premise or normality hypothesis is required.
\end{theorem}
\begin{proof}
First compare the two actual ambient-away seminorms. Every element of the
homogeneous-away algebra has an original homogeneous numerator $G/F^m$.
The preceding numerator identity and the original chart evaluation formula
give the same value $|G|/|F|^m$ on both sides. Extensionality therefore gives
equality of the whole affine seminorms. The actual native-away chart theorem
identifies the original normalized projective point with the image of that
same affine seminorm point, including its canonical residue valuation.
Apply the whole original inclusion square. The established projective
analytic/native comparison commutes with the original closed inclusion, and
mapping native points through a closed immersion is injective. Cancelling
that map proves the claimed whole-point identity.
\end{proof}
